% TOP_PROBLEM: {"id": 2304025, "title": "Research Problems in Function Theory — Problem 4.25", "category_id": 8, "category": {"id": 8, "name": "analysis", "display_name": "Analysis", "description": "Limits, continuity, calculus, and function theory.", "slug": "analysis", "order_index": 8, "created_at": "2026-07-31T15:26:25.671Z"}, "status": "open", "problem_number": "AMR-022-4025", "set_id": 13}
% FIRST_POSTED: 2026-10-01
% ENTRY_KIND: statement_only
% UnsolvedMath ID 2304025; source catalogue code AMR-022-4025
% !TeX program = lualatex
% Definition and sources only; this file is not an LLM solution attempt.
\documentclass[11pt,a4paper]{article}
\usepackage[margin=25mm]{geometry}
\usepackage{fontspec}
\setmainfont{lmroman10-regular.otf}[BoldFont=lmroman10-bold.otf,
 ItalicFont=lmroman10-italic.otf,BoldItalicFont=lmroman10-bolditalic.otf]
\usepackage{amsmath,amssymb,mathtools}
\usepackage{unicode-math}
\setmathfont{latinmodern-math.otf}
\AtBeginDocument{\let\setminus\smallsetminus}
\usepackage{xurl,hyperref}
\hypersetup{colorlinks=true,urlcolor=blue}
\setcounter{secnumdepth}{0}
\setlength{\parindent}{0pt}
\setlength{\parskip}{5pt}
\setlength{\emergencystretch}{3em}
\begin{document}
\section*{Research Problems in Function Theory — Problem 4.25}\label{2304025}
{\small UnsolvedMath ID 2304025 · AMR-022-4025 · Analysis · AMR Open Problem Lists}
\subsection{Definitions and mathematical statement}
Fix a real parameter $\lambda>0$. Let $\mathcal M$ be the collection
of all monic polynomials with integer coefficients, in any finite
degree. The constant polynomial $1$ is included, while the zero
polynomial is excluded. For $P\in\mathcal M$, define
\[
 J_\lambda(P)=\int_{-\pi}^{\pi}
       |1-e^{i\theta}|^{2\lambda}|P(e^{i\theta})|^2\,d\theta,
 \qquad I(\lambda)=\inf_{P\in\mathcal M}J_\lambda(P).
\]
The integral is unnormalized Lebesgue integration in the angle
variable; thus there is no factor $1/(2\pi)$. The weight is
nonnegative, vanishes only at $\theta=0$ modulo $2\pi$, and is
integrable for the stated range of $\lambda$.

Determine $I(\lambda)$ as a function of $\lambda$. A sharp answer
must give a lower bound valid simultaneously for all degrees and
all integer coefficient choices, together with polynomials whose
integrals approach that bound. Determine also whether the infimum
is attained, and describe extremal polynomials when they exist.
The degree is allowed to grow along an approximating sequence;
minimizing separately in each fixed degree is only an intermediate
step. No restriction is imposed on the locations of the roots,
on the signs of the integer coefficients, or on their magnitudes.
The monic condition is essential because arbitrary multiplication
by small constants would otherwise destroy the intended arithmetic
normalization.
\subsection{Short English statement}
Find the least weighted boundary square integral among all monic integer polynomials, with the degree unrestricted.
\subsection{Sources}
\begin{itemize}
\item[S1] ULAM.AI and the UnsolvedMath Contributors, supplied problems.json, record 2304025 (AMR-022-4025), AMR Open Problem Lists. Catalogue identity and source transcription; CC BY 4.0.
\url{https://huggingface.co/datasets/ulamai/UnsolvedMath}.
\item[S2] Hayman - Research Problems in Function Theory (2018), Problem 4.25. Original source indexed by AMR-022-4025.
\url{https://arxiv.org/abs/1809.07200}.
\end{itemize}
\end{document}
